\documentclass[10pt,journal,twocolumn]{IEEEtran}
\usepackage{fontspec}
\setmainfont[Path=fonts/,BoldFont=timesbd.ttf,ItalicFont=timesi.ttf,BoldItalicFont=timesbi.ttf]{times.ttf}
\usepackage{amsmath,amssymb,booktabs,array,graphicx}
\usepackage[letterpaper,margin=1in,columnsep=.25in]{geometry}
\usepackage{balance}
\raggedbottom
\usepackage[hidelinks,unicode]{hyperref}
\newcommand{\DN}{\,\mathrm{DN}}
\newcommand{\mean}{\operatorname{mean}}
\newcommand{\BCE}{\operatorname{BCEWithLogits}}
\newcommand{\R}{\mathbb{R}}
\newcommand{\figwide}[4]{\begin{figure*}[!t]\centering\includegraphics[width=#2\textwidth]{figures/#1.pdf}\caption{#3}\label{#4}\end{figure*}}
\newcommand{\figone}[3]{\begin{figure}[!t]\centering\includegraphics[width=\columnwidth]{figures/#1.pdf}\caption{#2}\label{#3}\end{figure}}
\setcounter{dbltopnumber}{2}
\renewcommand{\dbltopfraction}{0.92}
\renewcommand{\dblfloatpagefraction}{0.85}
\makeatletter
\setlength{\@dblfptop}{0pt}
\setlength{\@dblfpbot}{0pt}
\newenvironment{inlinetable}{\par\medskip\noindent\begin{minipage}{\columnwidth}\def\@captype{table}}{\end{minipage}\par\medskip}
\makeatother
\setlength{\emergencystretch}{1em}
\hyphenation{black-level photo-electron photo-electrons}


\usepackage{xeCJK}
\setCJKmainfont[Path=fonts/,BoldFont=msyh.ttc]{simsun.ttc}
\renewcommand{\figurename}{图}
\renewcommand{\tablename}{表}
\hypersetup{pdftitle={基于原生像素证据的时域稳定 RAW 黑电平估计},pdfauthor={姜尧耕},pdfsubject={黑电平估计与时域 RAW 校正}}
\begin{document}
\title{{\small\ttfamily repo/update: \href{https://github.com/y-g-jiang/Video-RAW-Black-Level-Correction-Net}{https://github.com/y-g-jiang/}\\\href{https://github.com/y-g-jiang/Video-RAW-Black-Level-Correction-Net}{Video-RAW-Black-Level-Correction-Net}}\\[0.8em]基于原生像素证据的时域稳定 RAW 黑电平估计}
\author{\href{https://y-g-jiang.github.io}{姜尧耕}\\\url{https://y-g-jiang.github.io}}
\maketitle
\begin{abstract}
亚 DN 级黑电平误差经曝光提亮后会引起暗部偏色，并在 RAW 视频中产生时域波动。本文结合多尺度空间特征与带符号 Bayer 观测的直接注意力检索，以传感器原生数码值预测四个 CFA 通道偏移。查询在帧间交换信息并重新读取各帧观测，在共享权重下支持独立单帧（T1）与联合七帧（T7）推断。在单场景、14 组光子转移曲线与 13 种偏移模式组成的测试矩阵上，预训练模型将非零偏移 MAE 从 0.649947 DN 降至 0.054768 DN。三台相机的 +6 EV 曝光阶梯实例相对对齐的同场景 0 EV 参考图，暗部色度距离降低 16.54\%--36.09\%。在八场景开发集的 105k 工作点，T7 将相邻变化误差从 0.073666 DN 降至 0.054778 DN，绝对误差从 0.108795 DN 降至 0.102115 DN。在 Z63 衍生连续序列上，五类七帧 T1 后处理滤波器采用留一场景选参与响应校准。长序列配对扰动在与评分相同的滑动窗中心测量响应，汇总逐 CFA 增益与 T7 的最大差为 0.0353。T7 的绝对／相邻变化误差为 0.1270/0.0480 DN，滤波器为 0.1328--0.1336/0.0543--0.0587 DN。这些结果量化了共享权重估计器的精度、时域响应与计算成本。
\end{abstract}
\begin{IEEEkeywords}
黑电平估计，计算成像，RAW 图像处理，低光照成像，时域稳定性，原生像素证据。
\end{IEEEkeywords}

\section{引言}
黑电平校准影响数字成像的光度准确性与暗部色彩。数字曝光提亮可将亚数码值（DN）误差放大为可见色偏~\cite{ptc_quant}，视频中逐帧独立估计还可能使校正量产生时域波动。本文直接在传感器原生单位下估计四个彩色滤光阵列（CFA）通道的残差偏移，并在去马赛克与非线性 ISP 渲染前完成减偏移校正。

BLBE 在深度 RAW 去噪前估计四个 CFA 通道偏移，以减小低光色偏~\cite{blbe}。RAW 恢复与物理噪声模型研究单幅图像恢复~\cite{sid,sidd,eld,noiseflow,unprocessing}，连拍与视频方法进一步利用时域观测重建图像内容~\cite{smd,fastdvd,crvd}。本文利用多帧观测估计每帧四个标量校正量：SimpleGate 空间编码器~\cite{naf}提供多尺度场景信息，受集合聚合~\cite{set,perceiver}启发的紧凑查询直接读取原生分辨率的带符号 Bayer 单元。

在视频推断中，查询在帧间交换信息，再重新读取各帧原生观测。独立单帧（T1）与联合七帧（T7）推断共享同一套权重和空间编码器。实验分别在 PTC 合成矩阵、三相机曝光阶梯实例与 Z63 衍生开发序列上评价偏移精度、暗部色彩及时间误差和响应，并以固定 T1 预测的经典七帧后处理作为时域比较。

\figwide{architecture}{.98}{模型整体架构：空间特征与原生 Bayer 单元观测共同支持每帧 8 个紧凑的可学习查询向量。T7 模式在 7 帧窗口内交换表征并重新检索各帧原生空间特征，T1 模式各帧独立推断。两种模式完全共享同一套权重，并在非线性渲染前从带符号 RAW 中减去预测的 4 通道标量偏移。}{fig:architecture}

\section{观测模型与估计目标}
\subsection{原生数值域下的四通道校正模型}
设 $y_{t,c}(p)$ 为视频第 $t$ 帧、通道 $c\in\{R,G_1,G_2,B\}$、通道内空间坐标 $p$ 处减去静态参考值后的带符号 RAW 观测值。传感器观测表示为
\begin{equation}
 y_{t,c}(p)=x_{t,c}(p)+\delta_{t,c}+\epsilon_{t,c}(p),
 \label{eq:observation}
\end{equation}
其中 $x$ 为以 DN 为单位的场景无噪信号，$\delta$ 为待估计的残差偏移，$\epsilon$ 代表包含读出噪声、光子散粒噪声与量化误差的复合噪声。

给定估计偏移向量 $\hat\delta$，校正后的 RAW 帧通过逐通道减去标量常数获得
\begin{equation}
 z_{t,c}(p)=y_{t,c}(p)-\hat\delta_{t,c}.
 \label{eq:correction}
\end{equation}
因此，对于同一帧、同一通道内的任意两个空间位置 $p$ 与 $q$，均恒等满足
\begin{equation}
 z_{t,c}(p)-z_{t,c}(q)=y_{t,c}(p)-y_{t,c}(q).
 \label{eq:invariance}
\end{equation}
在下游非线性截断或量化之前，这种加性常数平移在严格保全通道内部空间对比度与纹理细节的同时，准确重置了跨通道的色彩中性基准。

\subsection{绝对精度、时域稳定性与动态响应}
微弱的场景光照与正向残差偏移在零点附近具有相似的数值表现。分离二者依赖于空间纹理结构、极低信号采样点、信噪比统计特性以及跨帧时域上下文。由于各场景包含的信息丰度不同，评测必须涵盖非零动态轨迹、真实零偏移检验以及受控扰动响应。

设 $e_{t,c}=\hat\delta_{t,c}-\delta_{t,c}$ 为估计残差。对于序列集合 $n$，我们定义绝对误差 $A$、时域均值误差 $C$、相邻帧变动误差 $J$ 以及输出波动 $S$
\begin{align}
 A&=\mean_{n,t,c}|e_{n,t,c}|,\label{eq:A}\\
 C&=\mean_{n,c}\left|\mean_t e_{n,t,c}\right|,\label{eq:C}\\
 J&=\mean_{n,t,c}|e_{n,t+1,c}-e_{n,t,c}|,\label{eq:J}\\
 S&=\mean_{n,t,c}|\hat\delta_{n,t+1,c}-\hat\delta_{n,t,c}|.\label{eq:S}
\end{align}
其中 $A$ 衡量整体绝对准确度，$C$ 反映长时静态偏差，$J$ 衡量相邻帧偏移变化的估计误差，$S$ 衡量原始输出的绝对变动幅度。此外，成对阶跃扰动增益用于量化模型对突发动态偏移的跟随能力，而在真实零偏移测试下的 MAE 及其 95 分位数则用于检验不存在残差时的误校正风险。

\figwide{video_demo_z63_empirical_zh_v12}{.98}{Z63 暗场轨迹衍生序列可视化：选取经验变化区间内的三个代表性固定时点（第 255、295、335 帧），在相同传感器空间坐标下展示辐射代理、带符号含噪 RAW 以及注入真值校正后的结果。场景中位数 1 DN，显示增益 +3 EV。底部曲线给出该区间全部 120 帧的偏移轨迹（29.97 帧/秒）。八个开发序列上的模型估计效果对比见图~\ref{fig:z63_empirical}。}{fig:video_demo}

\section{结合空间与原生像素证据的估计架构}
\label{sec:method}

\subsection{多尺度空间上下文与原生像素记忆}
四阶段 SimpleGate 编码器~\cite{naf}提取下采样 8 倍的高维特征网格。与此同时，原生通道保留每一个 $2\times2$ Bayer 四合一单元的带符号读数、归一化空间坐标及传感器元数据，构成 22 维局部描述符，经线性投影为 16 维键向量与 28 维值向量。对于查询向量 $q$ 与单元特征 $d_i$，原生读取定义为
\begin{equation}
\begin{aligned}
 R_n(q,N)&=W_o\!\left(\sum_{i=1}^{M}\alpha_i(q)v(d_i)\right),\\
 \alpha_i(q)&=\frac{\exp a_i(q)}{\sum_j\exp a_j(q)}.
\end{aligned}\label{eq:read}
\end{equation}
其中 $a_i$ 为查询与键之间的注意力对数几率。Softmax 归一化跨越整幅裁切窗口内的全部 $M$ 个 Bayer 单元，采用数值稳定的分块 LogSumExp 累加实现。空间特征读取 $R_g$ 同步提供大感受野语义上下文。两条通路在代表性局部视窗内完成联合推断，随后将所得标量常数广播应用于整幅全分辨率 RAW 图像。

\subsection{时域表征交换与记忆重读机制}
每帧由 8 个学习得到的 192 维全局查询通过检索空间网格与原生像素记忆完成初始化
\begin{equation}
 Q_t^0=R_g(Q_*,G_t)+0.1R_n(Q_*,N_t).
 \label{eq:initialread}
\end{equation}
随后，时域交换模块 $E_r$ 借助四头自注意力机制以及时域位置编码与段标识更新各帧查询表征。在前三个交换模块之后，更新后的时域查询均会返回重新检索本帧的空间与原生记忆
\begin{align}
 V_t^r&=E_r(Q^{r-1})_t,\nonumber\\
 Q_t^r&=V_t^r+0.25\big[R_g(V_t^r,G_t)
                 +0.1R_n(V_t^r,N_t)\big].
 \label{eq:reread}
\end{align}
第四个交换模块之后，最终预测头综合更新后的查询表征与空间网格的均值、标准差。邻帧信息参与查询更新，每次注意力重读均访问本帧的原生空间证据。

在单帧模式（T1）下，每帧被赋予独立的时域段掩码，而在七帧联合模式（T7）下，连续七帧被编入同一个共享时域段。两类配置完全复用相同的交换计算结构。因此，对多帧批次执行 T1 推断在数学上严格等价于逐帧独立运行，邻帧数据的任何扰动均不会穿透到当前帧的 T1 估计值中。在训练阶段，中间交换层还接入辅助预测头以提供深度监督。

\subsection{有界输出映射与残差监督损失}
设模型输出 logits 为 $s$，输出边界为 $B>0$，预测头与训练损失函数设计为
\begin{align}
 \hat\delta&=B\tanh(s/B),\label{eq:head}\\
 \ell_B(s,\delta)&=B^2\BCE\!\left(\frac{2s}{B},\frac{\delta/B+1}{2}\right).
 \label{eq:loss}
\end{align}
将损失函数对网络输出 $s$ 求导可得
\begin{equation}
 \frac{\partial\ell_B}{\partial s}=\hat\delta-\delta.
 \label{eq:gradient}
\end{equation}
该映射将输出严格约束在 $[-B,B]\DN$ 区间内，其损失梯度大小精确等于估计值与真实值之间的残差。

预训练单帧模型采用 $B=1\DN$ 的输出边界，在时域适应阶段扩展为 $B=3\DN$，以容纳累积常数偏移与动态时域漂移。表\ref{tab:metrics}中的基线指标统一在 $B=3\DN$ 边界下测定。

\figwide{real_color}{.98}{S5 II、OM-1 与 S9 曝光阶梯实例。各列依次为同场景 0 EV 参考图、采用元数据黑电平的六档欠曝输入经 +6 EV 恢复后的结果，以及相同显影条件下的预训练 T1 校正结果。}{fig:color}

\section{时域适应与实验设计}
\label{sec:adapt}

\subsection{目标数值域的合成数据构建}
时域适应训练采用 48 个训练场景与 8 个互不重叠的开发验证场景。线性高动态 RGB 辐射代理依次经过相机平移、旋转、缩放、双轴各 2 倍空间降采样积分以及 Bayer 采样，生成 $320\times480$ 像素的 RAW 视窗，重点刻画准静态构图与温和运镜。

传感器噪声与偏移目标在 Bayer 采样后注入。以尼康 Z6 III 相机工作点为例，静态黑白电平分别为 1008 与 16383 DN，满阱电荷量为 90,000 电子，对应转换增益 $k=0.170833\DN$/电子。高斯读出噪声标准差按序列在 $U[1.2,1.8]\DN$ 内随机采样。通道偏移构建为 $\delta_{t,c}=m_c+r_{t,c}$，其中 $m_c$ 为幅度属于 $U[0,1]\DN$ 的序列恒定基线，而 $r_{t,c}$ 来源于实机暗场记录提取的动态四通道耦合漂移。场景中位亮度覆盖 0.25 至 64 DN。泊松散粒噪声与高斯读出噪声在整数截断前施加，完整保留负值读数。详细采样参数分布与轨迹生成流程见补充材料。

\subsection{模型训练与检查点指标分析}
时域适应过程加载预训练 T1 权重，将输出范围扩展为 $B=3\DN$，并保留主干特征网络与 Adam 优化器状态。三条监督通路作用于相同的 7 帧 4 通道真值
\begin{equation}
 \mathcal L=0.25\mathcal L_{\mathrm{frame}}
          +0.375\mathcal L_{\mathrm{T1}}
          +0.375\mathcal L_{\mathrm{T7}}.
 \label{eq:training}
\end{equation}
模型包含 232 万个可训练参数。经过开发评估，本文将 105k 检查点固定为时域分析工作点。表\ref{tab:metrics}同时对比了按基线误差容限准则选出的早期 46k 模型，补充材料图 S2 进一步比较了 102k、105k 与 106k 工作点的误差与响应数值。

\figwide{ptc_matrix}{.98}{预训练 T1 估计器在固定测试场景下遍历 14 组 PTC 噪声参数与 13 种有符号偏移模式的四通道平均绝对误差（MAE）矩阵。色彩刻度统一定为 0--1 DN。}{fig:ptc}
\figwide{temporal}{.98}{八个开发验证场景在 105k 检查点下的时域性能分布：(a) 场景级相邻帧变动误差 $J$，(b) 原始输出变动 $S$，(c) 场景中位亮度为 2 DN 时的中心化相对响应增益。}{fig:temporal}

\subsection{评测协议}
PTC 矩阵由 OM-1 单场景、14 组噪声参数与 13 种偏移模式构成，包括 168 组非零与 14 组真实零偏移条件。短窗时域评测包含八个场景、每场景两条录制轨迹的 16 个七帧片段。连续评测为每场景一段 Z63 衍生 120 帧序列，评分中心为第 6--113 帧。各方法共用 Bayer 观测与三帧前视。五类 T1 后处理包括均值、中值、对称指数、随机游走 Kalman 和跳变混合状态空间平滑器，时间支持均为七帧。滤波参数与逐 CFA 响应校准由七个场景拟合，再应用于余下场景。该留一场景评测在八个开发场景上进行，网络固定为 105k 检查点。实拍评价采用 S5 II、OM-1 与 S9 曝光阶梯。 120 帧配对扰动序列在相同滑动窗中心测量网络与滤波器响应。

\figwide{z63_empirical_t1_t7_zh}{.98}{105k 在 Z63 衍生连续开发序列上的结果。左：固定第一场景的四 CFA 真值及 R 通道跟踪，偏移纵轴统一为 $-2$ 至 $2$ DN。右：T1、T7 与五类响应校准后处理的八场景相邻变化误差。每个场景的参数由另外七个场景确定，各方法共用观测与 108 个有效中心帧。}{fig:z63_empirical}

\section{实验与结果分析}
\label{sec:results}

\subsection{实拍低光相机暗部校正验证}
预训练单帧模型在 S5 II、S9 的 ISO 100 及 OM-1 的 ISO 200 六档欠曝 RAW 上评价。每台相机提供一个棚拍场景及其正常曝光的 0 EV 图像，欠曝输入通过 +6 EV 线性增益恢复。0 EV 场景图用于配准、亮度匹配与暗部色度距离计算。

图\ref{fig:color}显示四通道校正后暗部偏色减轻。相对对齐的同场景 0 EV 参考图，S5 II 的 CAM16-UCS 暗部色度距离从 4.811925 降至 4.015822，OM-1 从 14.256274 降至 10.228349，S9 从 12.079137 降至 7.720369（表\ref{tab:color}），在固定显影与匹配条件下分别下降 16.54\%、28.25\% 和 36.09\%。
\begin{table}[!t]
\centering
\caption{+6 EV 曝光恢复的暗部色度距离（CAM16-UCS）。每台相机采用同场景 0-EV 图像作为参考。}\label{tab:color}
\footnotesize
\setlength{\tabcolsep}{4pt}
\begin{tabular}{lrrr}
\toprule
相机 & 校正前 & 校正后 & 降幅 \\
\midrule
S5 II & 4.811925 & 4.015822 & 16.54\% \\
OM-1 & 14.256274 & 10.228349 & 28.25\% \\
S9 & 12.079137 & 7.720369 & 36.09\% \\
\bottomrule
\end{tabular}
\end{table}



\subsection{合成 PTC 矩阵基准评测}
在覆盖全面的 PTC 合成矩阵测试中（图~\ref{fig:ptc}），预训练 T1 估计器在全部 168 组非零偏移条件下均实现了大幅精度改善，将总体平均绝对误差从 0.649947 DN 降低至 0.054768 DN。在 14 组真实零偏移测试条件下，平均误校正误差为 0.006281 DN，整张参数网格上的最大单元格平均误差为 0.162939 DN。


\subsection{时域精度与动态响应}
\begin{table}[!t]
\centering
\caption{开发集上同权重 T1 与 T7 的误差统计（DN）。其中 $A$、$C$、$J$ 评估经验轨迹，零偏移两列评估独立的真零输入，适配前输出界限为 3 DN。}\label{tab:metrics}
\footnotesize
\setlength{\tabcolsep}{2pt}
\begin{tabular}{lrrrrr}
\toprule
模式 & $A$ & $C$ & $J$ & 零 MAE & 零 P95 \\
\midrule
适配前 T1 & 0.2928 & 0.2575 & 0.1734 & 0.0812 & 0.2011 \\
适配前 T7 & 0.2920 & 0.2566 & 0.1709 & 0.0800 & 0.1858 \\
46k T1 & 0.0970 & 0.0831 & 0.0814 & 0.0163 & 0.0680 \\
46k T7 & 0.0954 & 0.0852 & 0.0713 & 0.0161 & 0.0773 \\
105k T1 & 0.1088 & 0.0943 & 0.0737 & 0.0084 & 0.0295 \\
105k T7 & 0.1021 & 0.0925 & 0.0548 & 0.0081 & 0.0278 \\
\bottomrule
\end{tabular}
\end{table}


在 105,000 次更新检查点下，模型从单帧独立 T1 推断切换为联合七帧 T7 推断后，相邻帧变动误差 $J$ 由 0.073666 DN 降至 0.054778 DN（降幅 25.64\%），绝对误差 $A$ 由 0.108795 DN 降至 0.102115 DN（降幅 6.14\%）。配对场景 Bootstrap 95\% 区间显示，T7 相对 T1 的差异在 $J$ 指标上落在 $[-0.025676,-0.011593]\DN$，在 $A$ 指标上落在 $[-0.010518,-0.003028]\DN$。全部 8 个场景的 $J$ 指标均呈现改善，其中 7 个场景的 $A$ 指标同步改善。作为对比，较早的 46,000 次更新模型呈现 $A=0.095354\DN$ 与 $J=0.071338\DN$（详见表~\ref{tab:metrics}）。

在 2 DN 弱信号场景下，105k T7 相对响应增益在四个 CFA 通道上分布于 0.505 至 0.608，低于 T1 的 0.785 至 0.850，四通道同向偏移增益为 0.745--0.763，低于 T1 的 0.797--0.831。在 64 DN 较亮场景下，T7 同向增益衰减至 0.076--0.162。在偏移幅度位于 $(0,0.1]\DN$ 的微小扰动区间内（包含 13 个相关帧通道），T7 的 MAE 为 0.109206 DN，其中有 7 组预测误差大于零校正基准。


\subsection{基于实测暗场轨迹的连续序列评测}
在八段 Z63 衍生连续序列上，T1 与 T7 的 $A$ 分别为 0.132780/0.126973 DN，$J$ 为 0.060524/0.048024 DN，固定七帧均值的 $A/J$ 为 0.163986/0.115446 DN。表\ref{tab:filters}比较以配对滑窗中心响应完成校准的五类滤波器，其汇总共同与相对增益在各 CFA 上与 T7 的最大差为 0.0353，$A$ 为 0.132834--0.133597 DN，$J$ 为 0.054323--0.058703 DN。校准跳变平滑器的 $A/J$ 为 0.132834/0.054323 DN，T7 在七个场景降低 $A$、六个场景降低 $J$。T7 减去该对照的配对场景 Bootstrap 95\% 区间为：$A$ 的 $[-0.008743,-0.003141]\DN$、$J$ 的 $[-0.010941,-0.002215]\DN$。相对其余四类校准滤波器，T7 在全部八个场景降低 $J$。
\begin{table}[!t]
\centering
\caption{七帧响应校准对照的留一场景开发集结果。$A,J$ 单位为 DN，滑窗中心扰动增益取场景与四 CFA 平均，耗时为流式每中心帧毫秒数。}\label{tab:filters}
\footnotesize
\setlength{\tabcolsep}{2.6pt}
\begin{tabular}{lrrrrr}
\toprule
方法 & $A$ & $J$ & $\bar g_r$ & $\bar g_c$ & ms \\
\midrule
T1 & 0.1328 & 0.0605 & 0.781 & 0.549 & 11.76 \\
T1 + 均值 & 0.1331 & 0.0586 & 0.745 & 0.553 & 11.76 \\
T1 + 中值 & 0.1336 & 0.0587 & 0.745 & 0.553 & 11.76 \\
T1 + 指数 & 0.1331 & 0.0584 & 0.745 & 0.553 & 11.76 \\
T1 + Kalman & 0.1331 & 0.0584 & 0.745 & 0.553 & 11.76 \\
T1 + 跳变模型 & 0.1328 & 0.0543 & 0.753 & 0.552 & 11.77 \\
T7 & 0.1270 & 0.0480 & 0.754 & 0.552 & 41.46 \\
\bottomrule
\end{tabular}
\end{table}


在补充图 S1 中，我们进一步评估了零偏移、静态背景 $m$、实测漂移 $m+r$ 以及叠加六次指定通道阶跃干预的组合序列 $m+r+j$。在相同的 7 帧时间窗口内，当序列仅包含完全静态的背景偏移时，T1 滑动平均能够有效平抑读出噪声，获得更低的 $J$（0.008207 DN 对比 T7 的 0.040470 DN），但在包含六次指定阶跃干预的序列（$m+r+j$）中，T7 的 $A/J$ 达到 0.155075/0.046149 DN，优于滑动平均的 0.163821/0.080308 DN。

\section{结论}
本文结合空间上下文与带符号 Bayer 观测检索，以传感器原生单位估计四通道偏移，并通过共享查询支持独立单帧与联合七帧推断。单场景 PTC 测试量化偏移估计精度，三台相机曝光阶梯实例展示暗部色度误差下降。在 Z63 衍生开发序列上，T7 的平均绝对误差与相邻变化误差低于五类选参并校准响应的 T1 后处理滤波器。配对滑窗中心扰动同步量化共同响应与经验轨迹响应，刻画联合推断的精度、响应与计算成本关系。

\balance
\begin{thebibliography}{14}
\bibitem{ptc_quant} Y. Jiang, ``Impact of ISP affine re-quantization on photon transfer characteristics and perceptual quality,'' Zenodo, Dec. 26, 2025, doi: 10.5281/zenodo.18949317.
\bibitem{blbe} M. Mohammadi \emph{et al.}, ``Why low-light cameras go color blind: Removing color bias in raw denoising,'' 2026, arXiv:2607.11090v1. [Online]. Available: \url{https://arxiv.org/abs/2607.11090}
\bibitem{sid} C. Chen, Q. Chen, J. Xu, and V. Koltun, ``Learning to see in the dark,'' in \emph{Proc. IEEE/CVF Conf. Computer Vision and Pattern Recognition (CVPR)}, 2018, pp. 3291--3300.
\bibitem{sidd} A. Abdelhamed, S. Lin, and M. S. Brown, ``A high-quality denoising dataset for smartphone cameras,'' in \emph{Proc. IEEE/CVF CVPR}, 2018, pp. 1692--1700.
\bibitem{eld} K. Wei, Y. Fu, J. Yang, and H. Huang, ``A physics-based noise formation model for extreme low-light raw denoising,'' in \emph{Proc. IEEE/CVF CVPR}, 2020, pp. 2758--2767.
\bibitem{noiseflow} A. Abdelhamed, M. A. Brubaker, and M. S. Brown, ``Noise Flow: Noise modeling with conditional normalizing flows,'' in \emph{Proc. IEEE/CVF International Conf. Computer Vision (ICCV)}, 2019, pp. 3165--3173.
\bibitem{unprocessing} T. Brooks, B. Mildenhall, T. Xue, J. Chen, D. Sharlet, and J. T. Barron, ``Unprocessing images for learned raw denoising,'' in \emph{Proc. IEEE/CVF CVPR}, 2019, pp. 11036--11045.
\bibitem{smd} C. Chen, Q. Chen, M. N. Do, and V. Koltun, ``Seeing motion in the dark,'' in \emph{Proc. IEEE/CVF ICCV}, 2019, pp. 3185--3194.
\bibitem{fastdvd} M. Tassano, J. Delon, and T. Veit, ``FastDVDnet: Towards real-time deep video denoising without flow estimation,'' in \emph{Proc. IEEE/CVF CVPR}, 2020, pp. 1354--1363.
\bibitem{crvd} H. Yue, C. Cao, L. Liao, R. Chu, and J. Yang, ``Supervised raw video denoising with a benchmark dataset on dynamic scenes,'' in \emph{Proc. IEEE/CVF CVPR}, 2020, pp. 2301--2310.
\bibitem{naf} L. Chen, X. Chu, X. Zhang, and J. Sun, ``Simple baselines for image restoration,'' in \emph{Proc. European Conf. Computer Vision (ECCV)}, 2022, pp. 17--33.
\bibitem{set} J. Lee, Y. Lee, J. Kim, A. R. Kosiorek, S. Choi, and Y. W. Teh, ``Set Transformer: A framework for attention-based permutation-invariant neural networks,'' in \emph{Proc. International Conf. Machine Learning (ICML)}, vol. 97, 2019, pp. 3744--3753.
\bibitem{perceiver} A. Jaegle, F. Gimeno, A. Brock, O. Vinyals, A. Zisserman, and J. Carreira, ``Perceiver: General perception with iterative attention,'' in \emph{Proc. ICML}, vol. 139, 2021, pp. 4651--4664.
\bibitem{cam16} C. Li \emph{et al.}, ``Comprehensive color solutions: CAM16, CAT16, and CAM16-UCS,'' \emph{Color Research \& Application}, vol. 42, no. 6, pp. 703--718, 2017, doi: 10.1002/col.22131.
\end{thebibliography}

\end{document}
